%%%%%%%%%%%%%%%%%%%%%%%%%%%%%%%%%%%%%%%%%%%%%%%%%%%%%%%%%%%%%%%%
% [메모] 8/30 압축판 (COMPRESSION_RULES.md, 목표 1,500 단어).  구조는 v2
%   5소절 그대로: A Setup / B Volume / C Density (main) / D Pose & Path /
%   E Asset.  수치 출처 불변: RESULTS.md (E1~E10), gt.json, frozen vision 열,
%   figures/v2/*.json.  없는 값은 \Pending / \rev.  실물 로봇 데이터 없음.
%   Exp E 는 시뮬-시뮬 대리 검증으로만 서술.  라벨 호환: sec:vol / sec:sim /
%   sec:pred 유지.  Table I 은 실제 표로 복원 (tables/table1_objects_v2.tex,
%   tab:objects 라벨은 표 안).  실물 비교 자리표시자 그림(구 Fig. 8)은 Fig. 7
%   캡션에 흡수 (fig:assetreal 라벨은 그림 안에 보존).  표 캡션·각주가 이미
%   말하는 문장은 산문에서 지웠고, 잘린 산문은 CUTS.md 에 보관.

\section{Experiments}
\label{sec:exp}

\subsection{Setup}
\label{sec:setup}

\begin{figure}[!htb]
\centering
\includegraphics[width=\linewidth]{fig3_setup.png}
\caption{Experimental platform and the four evaluation objects of
Table~\ref{tab:objects}. \rev{[Placeholder: the photo awaits the hardware
session.]}}
% [번역] 그림 3 — 실험 플랫폼과 평가 물체 4종(표 I). 실물 사진은 하드웨어
%   세션 대기 자리표시자. 물체 사양(링크/P/rank/R_min/참값 출처)과 힌지
%   각주는 표 I 로 복원.
\label{fig:setup}
\end{figure}

% Table 1 — 실험 물체 사양.  출처: gt.json (gt_list.md 를 정본으로 재계산),
%   rank / R_min 은 my_work/study_theory.py 의 수치 검증 결과.
\begin{table}[t]
\caption{Evaluation objects: links, independently unknown densities $P$,
single-configuration rank $1+\min(3,P)$ (Prop.~\ref{prop:rank}), configuration
lower bound $R_{\min}=\lceil (P-1)/3\rceil$ (Prop.~\ref{prop:rounds}), and
ground-truth source.}
% [번역] 평가 물체: 링크 수, 독립 미지 밀도 P, 단일 형상 rank 1+min(3,P)
%   (명제 1), 형상 개수 하한 R_min = ceil((P-1)/3) (명제 3), 참값 출처.
\label{tab:objects}
\centering
\footnotesize
\setlength{\tabcolsep}{4pt}
\begin{tabular}{lccccl}
\toprule
Object & Links & $P$ & rank & $R_{\min}$ & Ground truth \\
\midrule
Custom 2-link  & 2 & 3$^{\mathrm h}$ & 3 & 1 & CAD $+$ material \\
Custom 3-link  & 3 & 5$^{\mathrm h}$ & 4 & 2 & CAD $+$ material \\
Stand lamp     & 3 & 3 & 3 & 1 & scan $+$ disassembly \\
Laptop$^{\dagger}$ & 2 & 2 & 3 & 1 & scale $+$ part split \\
\bottomrule
\end{tabular}
\vspace{2pt}
\parbox{\linewidth}{\footnotesize \textit{Rank and $R_{\min}$ follow from
geometry alone, before any wrench is measured. $^{\mathrm h}$ includes the friction hinge
as its own unknown ($41\,\mathrm{g}$ each). The hinge occupies $7.4\,\mathrm{cm^3}$ at
$5.51\,\mathrm{g/cm^3}$---small, largely occluded, and mass-dense, the case that makes
part-wise identification hard rather than the case that makes it easy.
$^{\dagger}$ \rev{Laptop mesh and URDF are in preparation; it contributes ground truth
and vision-baseline rows only.}}}
\end{table}



The platform: a six-axis RB5-850E arm, wrist force/torque sensor, Robotiq
2F-85 gripper, and one RealSense D456; joint angles are tracked on the
RGB-D stream---measured, not commanded (Fig.~\ref{fig:setup}). All
results below are Drake simulation, running the estimator, selection rule,
and stopping test of Sec.~\ref{sec:estimation} unchanged; \rev{the hardware runs
are in preparation, and every real-robot quantity is marked pending}.
% [번역] 장비: 6축 RB5-850E + 손목 F/T 센서 + Robotiq 2F-85 + RealSense D456
%   한 대. 관절각은 RGB-D 추적기의 측정값(명령값 아님)이다(그림 3). 아래 모든
%   결과는 III-E 의 추정기·선택 규칙·정지 판정을 그대로 돌린 시뮬레이션
%   (Drake)이다. 하드웨어 실행은 준비 중, 실물 수치는 전부 pending.

We evaluate four articulated objects (Fig.~\ref{fig:setup},
Table~\ref{tab:objects}): two custom 3D
prints with per-part densities we set, a disassembly-weighed stand lamp from
a 3DGS scan, and a laptop; the custom objects give their friction hinge its
own part---the harder general case. Density
identification covers three of the four: \rev{the laptop's mesh and URDF
assets are in preparation, so it contributes ground-truth and
vision-baseline rows only}.
% [번역] 평가 물체 4종(그림 3): 부위별 밀도를 지정한 커스텀 3D 프린트 2종,
%   분해 계량 참값의 3DGS 스캔 스탠드 램프, 노트북. 커스텀 둘은 마찰 힌지를
%   독립 part 로 둔다 — 작고 가려지고 질량이 몰린, 더 어려운 일반 경우. 밀도
%   식별은 4종 중 3종: 노트북 mesh/URDF 는 준비 중이라 참값·vision baseline
%   행만 기여한다.

Metrics: per-part relative density error; $95\%$ half-width and its
coverage (a confident wrong answer is worse than an uncertain one); rounds
to the target half-width $\tau=0.5\%\times P$; asset trajectory RMSE. All
methods share the same fixed part volumes, so density error and mass error
are the same number (Table~\ref{tab:main} footnotes). Four objects admit no
statistical-significance claim, so tables show per-part rows, not pooled
means.
% [번역] 지표: 부위별 상대 밀도오차, 보고 95 % 반폭과 경험적 커버리지(자신
%   있게 틀린 답이 불확실한 답보다 나쁘다), 목표 반폭 tau = 0.5 % x P 까지의
%   라운드 수, asset 은 궤적 RMSE. 모두 같은 고정 부피로 채점되므로 밀도오차
%   = 질량오차(1.7e-16 불변성은 표 III 각주). 물체 4개로는 통계적 유의성을
%   주장할 수 없어 표는 평균 아닌 부위별 행을 드러낸다.

Baselines are PUGS~\cite{shuai2025pugs}, SiPhy~\cite{le2026siphy}, and
PhysX-Omni~\cite{cao2026physxomni} from vision, plus three ablations:
\emph{uniform} (one density per assembly~\cite{pfaff2025scalable}),
\emph{single} (one fixed configuration, unlimited reorientations and
repeats~\cite{nadeau2023sum}), and \emph{WLS} (our loop, ignoring the
regressor's own angle error, Eq.~\eqref{eq:tls}). Vision outputs lacking
per-part densities are converted in whichever direction favors the baseline
(Table~\ref{tab:main} footnotes).
% [번역] baseline: vision 셋(PUGS, SiPhy, PhysX-Omni)과 ablation 셋 —
%   uniform 은 조립체에 밀도 하나, single 은 형상 하나 고정 + 무제한
%   reorientation·반복, WLS 는 우리 루프에서 회귀행렬 자체의 각도 오차(식
%   (7))만 무시. part 별 밀도가 없는 vision 출력은 매번 해당 baseline 에
%   유리한 방향으로 변환한다(표 III 각주).


\subsection{Part-Wise Volume Estimation}
\label{sec:vol}

\begin{figure*}[!t]
\centering
\includegraphics[width=0.74\textwidth]{fig_volume.png}
\caption{Injected part-volume error (eight seeds): density error tracks it
near one-to-one; part mass, scored with the volumes the estimator used,
drifts to $3$--$6\%$ at $\delta=30\%$---first-order invariant, not exactly
so.}
% [메모] Fig. 4 — 파일: figures/v2/fig_volume.png.
\label{fig:volume}
\end{figure*}


The CAD solids and displacement measurements are scoring references, never
inputs. \rev{The mesh refinement and per-part
scaling implementation is being delivered by a collaborator, so the measured
per-part volume-error table is deferred; no claim below depends on it.} What we can
already quantify is the downstream cost: we perturb the
volumes the estimator \emph{believes} by $\pm\delta$ with random signs while
measurements come from the true plant, scoring masses with the volumes
actually used ($\delta\le30\%$; Fig.~\ref{fig:volume}).
% [번역] CAD 솔리드·수중 배제 실측은 채점 기준이지 입력이 아니다.
%   refinement/scaling 구현은 협업자 인계 중이라 부피오차 실측 표는 유보,
%   아래 주장은 그에 의존하지 않는다. 지금 정량화할 수 있는 것은 하류 비용이다.
%   추정기가 '믿는' 부피를 ±delta 무작위 부호로 흔들되 측정은 참 plant 에서
%   나오고 질량은 실제 쓴 부피로 채점한다(delta <= 30 %, 그림 4).

Density error follows volume error near-exactly one-to-one; part mass is far
less sensitive but not exactly invariant---$3$--$6\%$ at $\delta=30\%$, the
MAP prior's cancellation being inexact (Remark~\ref{rem:volume} covers only
the least-squares core)---and the mesh-derived inertia tensor degrades as
$s^2$ ($19.1\%$ at $\delta=30\%$).
% [번역] 밀도오차는 부피오차를 거의 1:1 로 따르고, 부위 질량은 훨씬 덜
%   민감하나 엄밀히 불변은 아니며 — MAP 사전분포의 상쇄가 부정확해 delta
%   30 % 에서 3~6 % (Remark 의 불변성은 최소제곱 핵심만 다룬다) — mesh 유도
%   관성 텐서는 s^2 로 나빠진다(delta 30 % 에서 19.1 %).

Volume error also costs \emph{rounds}: the worst part drags the loop, median
rounds growing from $1$ to $7$ on the 2-link object by $\delta=30\%$, with
three of eight 3-link seeds failing to converge there. Refinement is what keeps the
interaction count bounded.
% [번역] 부피오차는 '라운드' 도 치른다. 최악 part 가 루프를 끌어 2링크에서
%   delta 30 % 까지 라운드 중앙값이 1 → 7 로 늘고, 3링크는 거기서 8 시드 중
%   3 이 미수렴. 상호작용 횟수를 유한하게 묶는 것이 바로 정제다.


% Table 2 (part별 부피 오차) 는 협업자의 refinement/scaling 실측이 도착할 때까지
% 전 칸이 Pending 이라 지면만 차지한다. 파일은 tables/table2_volume_v2.tex 에
% 그대로 두었으니, 자료가 오면 아래 한 줄의 주석만 풀면 된다.
% % Table 2 — part별 부피 오차.  refinement/scaling 구현과 실측 오차는 협업자 인계 대기.
%   빈 칸을 지어내지 않는다. 자료가 오면 \Pending 만 교체하면 된다.
\begin{table}[t]
\caption{Per-part volume error [\%] against the CAD solids (custom objects) and the
displacement measurements (stand lamp). The vision baselines are scored on the volume
they themselves recover---PUGS by integrating its Gaussians, PhysX-Omni on its generated
shape---so the table fixes how much of each method's density error is already present
before any density is assigned. \rev{The refinement and per-part scaling stage is being
delivered by a collaborator; its rows are pending.}}
\label{tab:volume}
\centering
\footnotesize
\setlength{\tabcolsep}{4pt}
\begin{tabular}{llccccc}
\toprule
& & \multicolumn{3}{c}{Baselines} & \multicolumn{2}{c}{Ours} \\
\cmidrule(lr){3-5}\cmidrule(lr){6-7}
Object & Part & Bbox & PUGS & \shortstack{PhysX-\\Omni} & \shortstack{raw\\mesh} & \shortstack{$+$refine\\$+$scale} \\
\midrule
\multirow{3}{*}{Stand lamp}
 & Base    & \Pending & \Pending & \Pending & \Pending & \Pending \\
 & Support & \Pending & \Pending & \Pending & \Pending & \Pending \\
 & Head    & \Pending & \Pending & \Pending & \Pending & \Pending \\
\midrule
\multirow{3}{*}{2-link}
 & Long    & \Pending & \Pending & \NA & \Pending & \Pending \\
 & Short   & \Pending & \Pending & \NA & \Pending & \Pending \\
 & Hinge   & \Pending & \Pending & \NA & \Pending & \Pending \\
\midrule
\multirow{4}{*}{3-link}
 & Link 0  & \Pending & \Pending & \NA & \Pending & \Pending \\
 & Link 1  & \Pending & \Pending & \NA & \Pending & \Pending \\
 & Link 2  & \Pending & \Pending & \NA & \Pending & \Pending \\
 & Hinges  & \Pending & \Pending & \NA & \Pending & \Pending \\
\bottomrule
\end{tabular}
\vspace{2pt}
\parbox{\linewidth}{\footnotesize \textit{The ablation columns isolate refinement from
scaling. \rev{Laptop rows are omitted until its asset is available.} An independent
reproducibility check of the reconstruction stage---two runs of the same scan under
identical settings---gave a scale spread of $0.24254$ vs.\ $0.24447\,\mathrm{m}$ per
unit and a volume spread of $29.9$ vs.\ $30.4\,\mathrm{L}$, i.e.\ a $\pm5\%$ noise floor
that no downstream density estimator can remove.}}
\end{table}


\subsection{Part-Wise Density Identification}
\label{sec:sim}
\label{sec:pred}

\begin{figure*}[!t]
\centering
\includegraphics[width=0.74\textwidth]{fig_conv.png}
\caption{Worst-part reported half-width per round (median and interquartile
range, eight seeds). Dashed: round bound $R_{\min}$ of
Prop.~\ref{prop:rounds}; dotted: target $\tau(P)$.}
% [번역] 라운드에 대한 부위 최악 보고 반폭(seed 8개 중앙값·사분위 범위).
%   파선은 명제 3의 라운드 하한 R_min, 점선은 목표 tau(P). ("형상만으로
%   계산"은 서론·표 I 각주·IV-C 산문, "하한 등호 달성"은 IV-C 산문에.)
\label{fig:conv}
\end{figure*}

% Table 3 (main) — per-part density identification.
%   GT V/M/rho: gt.json (권위 있는 정본).  vision 열: tables/table_main_v2.tex 의
%   frozen 값을 그대로 옮김 (재채점 금지).  uniform/single/WLS/ours:
%   my_work/figures/v2/main_real.json (seeds=8, rel=5%), part별 상대 밀도오차 중앙값.
%   coverage 블록: 같은 JSON 의 95% 구간 경험적 커버리지.
%   laptop: 자산 준비 중 → 8-페이지 압축에서 행 2개를 각주 한 문장으로 이전
%   (GT·vision 수치 전부 각주에 보존, 원행은 CUTS.md).
\begin{table*}[!t]
\caption{Part-wise density identification (simulation). Ground-truth $V$, $M$,
$\rho=M/V$; method columns give the relative per-part density error [\%] (median
over eight seeds for interaction methods), lowest per row in bold.}
% [번역] 부위별 밀도 식별(시뮬레이션). 참값 V, M, rho = M/V. 방법 열은 부위별
%   상대 밀도오차 [%](상호작용 방법은 8시드 중앙값), 행별 최저는 굵게.
%   (밀도오차=질량오차 1.7e-16 불변성과 커버리지 블록 설명은 각주로 이동.)
\label{tab:main}
\centering
\scriptsize
\setlength{\tabcolsep}{4pt}
\begin{tabular}{llrrrrrrrrrr}
\toprule
& & \multicolumn{3}{c}{Ground truth} & \multicolumn{3}{c}{Vision baselines} & \multicolumn{4}{c}{Interaction (median over 8 seeds)} \\
\cmidrule(lr){3-5}\cmidrule(lr){6-8}\cmidrule(lr){9-12}
Object & Part & $V$ [cm$^3$] & $M$ [g] & $\rho$ [g/cm$^3$] & PUGS$^{\mathrm S}$ & SiPhy & \shortstack{PhysX-\\Omni} & Uniform & Single & WLS & Ours \\
\midrule
\multirow{3}{*}{Stand lamp}
 & Base    & 273.00 & 399.00 & 1.462 & 164.1 & 38.3 & 69.1  & 7.33  & 0.12 & 0.12 & \textbf{0.01} \\
 & Support &  84.50 &  85.00 & 1.006 & 10.7  & 27.3 & 183.1 & 34.64 & 1.49 & 1.27 & \textbf{0.23} \\
 & Head    &  64.10 &  87.00 & 1.357 & 36.4  & 48.5 & 58.0  & 0.21  & 0.73 & 0.55 & \textbf{0.08} \\
\midrule
\multirow{3}{*}{2-link$^{*}$}
 & Long                    & 135.72 & 49.00 & 0.361 & 7887.4 & 593.4 & \NA & 49.14 & \textbf{0.18} & 0.21 & \textbf{0.18} \\
 & Short                   & 108.12 & 42.00 & 0.389 & 7323.6 & 544.5 & \NA & 25.43 & \textbf{0.06} & 0.10 & \textbf{0.06} \\
 & Hinge$^{\mathrm h}$     &   7.45 & 41.00 & 5.506 & 423.8  & 54.5  & \NA & 51.22 & \textbf{0.34} & 0.35 & \textbf{0.34} \\
\midrule
\multirow{5}{*}{3-link$^{*}$}
 & Link 0                  & 288.12 & 100.00 & 0.347 & 4312.5 & 349.6 & \NA & 43.32  & 0.11 & \textbf{0.01} & \textbf{0.01} \\
 & Link 1                  & 208.40 &  92.00 & 0.442 & 3369.1 & 253.5 & \NA & 110.52 & 1.14 & 0.31 & \textbf{0.06} \\
 & Link 2                  & 162.28 &  69.00 & 0.425 & 3501.8 & 267.0 & \NA & 321.04 & 1.58 & 0.48 & \textbf{0.10} \\
 & Hinge 1$^{\mathrm h}$   &   7.45 &  41.00 & 5.506 & \multirow{2}{*}{178.2$^{\ddagger}$} & \multirow{2}{*}{71.7$^{\ddagger}$} & \NA & 46.45 & 0.08 & \textbf{0.04} & \textbf{0.04} \\
 & Hinge 2$^{\mathrm h}$   &   7.45 &  41.00 & 5.506 &  &  & \NA & 46.45 & 0.08 & \textbf{0.05} & 0.06 \\
\midrule
\multicolumn{5}{l}{\emph{Coverage of the reported $95\%$ interval}} & & & & & & & \\
\multicolumn{5}{l}{\quad Stand lamp} & \NA & \NA & \NA & \NA & 24/24 & 4/24  & 22/24 \\
\multicolumn{5}{l}{\quad 2-link}     & \NA & \NA & \NA & \NA & 16/16 & 14/16 & 16/16 \\
\multicolumn{5}{l}{\quad 3-link}     & \NA & \NA & \NA & \NA & 24/24 & 15/24 & 24/24 \\
\bottomrule
\end{tabular}
\vspace{2pt}
\parbox{\textwidth}{\scriptsize \textit{All methods share the same fixed part
volumes, so density error and mass error are the same number (maximum observed
difference $1.7\times10^{-16}$) and a single column reports both. The bottom block
counts, over all part--seed pairs, how often each method's reported $95\%$ interval
covers the truth. $^{\mathrm h}$ hinge modelled as its own
unknown. $^{\mathrm S}$ Stand-lamp PUGS masses integrate its frozen SAM-partitioned
Gaussian volume and point-wise density, scored at the midpoint of its predicted mass
range; other PUGS rows and all SiPhy rows distribute each object-level prediction
using ground-truth volume fractions (Sec.~\ref{sec:setup}). PhysX-Omni is N/A where
its generated shape admits no canonical part mapping. $^{\ddagger}$ vision baselines
were scored on the two hinges jointly. $^{*}$ preliminary pre-ballast custom-object
masses and densities; CAD volumes are fixed. On the 2-link object Ours and Single
are byte-identical (Sec.~\ref{sec:sim}). Single's intervals
cover only because they are wide (Sec.~\ref{sec:sim}), and its convergence flag
never consults the stopping rule.
The fourth object, the laptop (display\,/\,base: $V=354.10/798.60\,
\mathrm{cm^3}$, $M=379.25/1137.75\,\mathrm{g}$, $\rho=1.071/1.425\,\mathrm{g/cm^3}$;
vision errors PUGS $76.7/30.1$, SiPhy $182.5/108.0$, PhysX-Omni $417.3/611.8\,\%$),
contributes ground-truth and vision-baseline data only: \rev{its mesh and URDF are
in preparation, so all its interaction columns are pending}.}}
\end{table*}



Each object runs the closed loop over eight seeds ($5\%$ joint-angle error,
$\tau=0.5\%\times P$); Table~\ref{tab:main} is the main result.
Prop.~\ref{prop:rounds} reads off geometry (Table~\ref{tab:objects}): the
3-link object's five unknowns exceed the rank-four ceiling, so a second
configuration is \emph{necessary} before any wrench is simulated.
% [번역] 각 물체는 각도오차 5 %, tau = 0.5 % x P 로 여덟 시드 폐루프를 돌며
%   표 III 이 메인 결과다. 명제 3은 형상만으로 읽힌다(표 I): 3링크는
%   미지수 다섯이 rank 4 천장을 넘어 렌치 시뮬 전에 이미 두 번째 형상이
%   '필요' 하다.

The loop meets that bound (one round on the 2-link object, two elsewhere;
all seeds converge). Uniform fails by orders of magnitude: the mass redistributions it misses
stay invisible to shape until a second articulation configuration
(Prop.~\ref{prop:null}).
% [번역] 루프는 그 하한대로 2링크 한 라운드, 나머지 두 라운드에 끝나고 전
%   시드 수렴한다. uniform 은 자릿수로 틀린다 — 놓친 질량 재분배는 두 번째
%   관절 형상 전까지 형상에 안 보인다(명제 2).

On the 2-link
object \emph{ours is single}: three unknowns fit under the rank-four
ceiling, and the two columns of Table~\ref{tab:main} are
byte-identical---the positive direction of Prop.~\ref{prop:rounds}, not a
win; only uniform is separated there. Single's failure appears on the 3-link
object, in the half-width rather than the point error: its worst-part error
is a respectable $1.58\%$, but its half-widths for the three links are
$2.6$, $28.7$, and $42.5\%$---one configuration cannot \emph{certify} five
unknowns.
% [번역] 2링크에서는 'ours 가 곧
%   single' — 미지수 셋이 rank 4 천장 아래라 표 III 두 열은 바이트 단위로
%   같다. 명제 3의 양의 방향이지 승리가 아니며 갈리는 것은 uniform 뿐.
%   single 의 실패는 3링크에서, 점오차 아닌 반폭에서다 — 최악 part 오차
%   1.58 % 는 준수하나 세 링크 반폭이 2.6 / 28.7 / 42.5 %: 형상 하나로
%   미지수 다섯을 '증명' 못 한다. (수렴 플래그가 무의미하다는 것과 single 의
%   커버리지가 폭 때문이라는 것은 표 III 각주가 싣는다.)

The strongest ablation evidence is the coverage block of
Table~\ref{tab:main}. WLS shares our loop, ignoring only the
errors-in-variables structure; its point errors stay within an order of
magnitude of ours, but its $95\%$ intervals miss the truth on up to $20$ of
$24$ part--seed pairs where ours cover at target width: the bias
Eq.~\eqref{eq:tls} removes does not shrink with rounds. Single's coverage
comes only from width. An estimator feeding a simulator must know when it
is wrong.
% [번역] 가장 강한 ablation 근거는 표 III 커버리지 블록이다. 루프를 공유하고
%   EIV 구조만 무시하는 WLS 는 점오차는 한 자릿수 안이지만 95 % 구간이
%   최대 24쌍 중 20쌍에서 참값을 놓친다(우리는 목표 폭으로 커버; 개별 수치는
%   표 III) — 식 (7) 이 제거하는 편향은 라운드를 쌓아도 안 준다. single 의
%   커버리지는 오직 폭 덕분이다. 시뮬레이터에 밀도를 넘길 추정기는 제가 틀린
%   때를 알아야 한다.

The hinge rows are the built-for case: densities recovered to $0.04$--$0.34\%$, intervals covering on all seeds, at
$3.0$--$3.8\%$ half-widths honestly reflecting the occluded part's weaker
observability. The stand lamp has no hinge body, so this rests on the custom
objects only.
% [번역] 힌지 행이 커스텀 물체를 만든 이유다. 밀도 0.04~0.34 %
%   복원, 구간은 전 시드에서 덮고, 반폭 3.0~3.8 % 는 가려진 part 의 약한
%   관측성을 정직하게 반영한다. 램프에는 힌지 body 가 없어 커스텀 둘에만
%   근거한다.

The stand lamp attains full rank in one configuration yet runs two: past
full rank, conditioning governs, set by joint-angle precision rather than
rank (Remark~\ref{rem:cond}, Fig.~\ref{fig:conv}).
% [번역] 램프는 형상 하나에서 full rank 인데 두 라운드를 돈다. full rank
%   이후는 조건수가 지배하고, 그것은 rank 아닌 각도 정밀도가 만든다
%   (Remark 3, 그림 5).

Synthetic chains of $P{=}2\ldots6$ unknowns separate the methods cleanly.
Through $P=3$ everything but uniform suffices; from $P=4$ the rank-four
ceiling of Cor.~\ref{cor:rigid} appears directly in \emph{single}, whose
worst-part error climbs $17.06\to32.02\to53.57\%$ as $P$ grows past four,
while WLS degrades $10.53\to11.79\to21.39\%$ and ours stays below $0.6\%$
throughout. At $P=4$ ours converges in a median of
$3$ rounds where offline D-optimal
design~\cite{gautier1992exciting,swevers1997optimal} needs $4$ and random
selection needs $21$ and fails once in eight; at $P=5$ the gap widens to $7$
against $10$ and random never converges ($0/8$). All three share the
criterion, so closed-loop re-selection is the only difference.
\rev{The $P=6$ arm is still computing.} Beyond $P=5$, though, no criterion---D, A, or E---reaches a $1\%$ target in
twelve rounds: at $\pm5\%$ angular error the binding resource is angular
precision, not the criterion.
% [번역] 미지수 P = 2..6 합성 사슬이 방법들을 깨끗하게 가른다. P <= 3 은 uniform
%   빼고 전부 충분하다. P = 4 부터 따름정리 2의 rank 4 천장이 single 에 그대로
%   나타나, 부위 최악 오차가 P 가 4를 넘어가며 17.06 -> 32.02 -> 53.57 % 로 오른다.
%   WLS 는 10.53 -> 11.79 -> 21.39 %, 우리는 전 구간 0.6 % 미만(0.25, 0.51, 0.59 %).
%   P = 4 에서 우리는 라운드 중앙값 3 으로 수렴하는데 고전 오프라인 D-최적 설계는 4,
%   무작위 선택은 21 이 필요하고 8회 중 1회는 실패한다. P = 5 에서는 격차가
%   7 대 10 으로 벌어지고 무작위는 예산 안에 아예 수렴하지 못한다(0/8).
%   셋이 기준을 공유하므로 차이는 폐루프 재선택 그 자체다. P = 6 arm 은 계산 중.
%   다만 기준 선택이 이 영역을 구제하지는 못한다 — P = 5 부터는 D, A, E 어느
%   기준도 12라운드 안에 목표 1% 에 닿지 못한다. 각도오차 ±5% 에서는 그만큼의
%   미지수를 가를 정보가 측정에 애초에 없기 때문이다. 구속하는 자원은 기준이
%   아니라 각도 정밀도다.



\subsection{Measurement Pose and Path Design}
\label{sec:posepath}

\begin{figure*}[!t]
\centering
\includegraphics[width=0.74\textwidth]{fig_path.png}
\caption{One row per object at its most informative colliding
configuration: joint-space paths, minimum clearance (dashed: straight;
solid: RRT-Connect), and end-link workspace trace.}
% [메모] Fig. 6 — 파일: figures/v2/fig_path.png.
\label{fig:path}
\end{figure*}

% Table 4 — 경로 비교 (figures/v2/path_compare.json 'all' 범위).
%   구 panel (a)(중력 방향 수확체감, study_tilt.py 실측)는 8-페이지 압축에서
%   삭제 — 수치는 IV-D 산문에 흡수, 원문은 CUTS.md.  E4b 주의 4건 반영:
%   관절한계 열은 두 계획기 모두 구조적으로 0 (볼록 상자) — 승리로 쓰지 말 것.
%   밀도오차 열: 두 계획기는 같은 측정 자세를 이으므로 추정치는 동일 —
%   각주로 설명 (숫자를 지어내지 않는다).
\begin{table}[!htb]
\caption{Straight-line joint interpolation vs.\ RRT-Connect over every pose pair
of the configuration grid (each cell: straight\,/\,RRT-Connect), with
per-configuration information gain $\Delta I$ and whitened-regressor condition
number $\kappa$.}
% [번역] 전체 형상 격자의 모든 자세쌍에서 직선 보간 대 RRT-Connect(칸마다
%   직선/RRT-Connect), 형상별 정보이득 Delta I 와 백색화 회귀행렬 조건수
%   kappa. (두 계획기의 밀도 추정 동일 문장은 각주로 이동.)
\label{tab:path}
\centering
\scriptsize
\setlength{\tabcolsep}{3pt}
\begin{tabular}{lccccc}
\toprule
Object & \shortstack{$\Delta I$ [nat] /\\cond.\ $\kappa$} & \shortstack{Coll.\\{[\%]}} & \shortstack{Pen.\\{[mm]}} & \shortstack{Clear.\\{[mm]}} & \shortstack{Time\\{[s]}} \\
\midrule
2-link & \shortstack{4.0--4.5\\57--324} & 25.0\,/\,\textbf{0.0} & 83.5\,/\,\textbf{0.0} & 0.6\,/\,13.1 & $\sim$0\,/\,0.12 \\
\addlinespace[2pt]
3-link & \shortstack{6.3--11.1\\$10^{16}$--$10^{18}$} & 20.9\,/\,\textbf{0.0} & 33.1\,/\,\textbf{0.0} & 11.3\,/\,14.1 & $\sim$0\,/\,0.11 \\
\addlinespace[2pt]
Stand lamp & \shortstack{6.9--9.3\\8.9--100} & 21.1\,/\,\textbf{0.0} & 93.8\,/\,\textbf{0.0} & $-0.8$\,/\,12.0 & $\sim$0\,/\,0.43 \\
\bottomrule
\end{tabular}
\vspace{2pt}
\parbox{\linewidth}{\footnotesize \textit{Coll.\ = fraction of pose pairs whose path
touches an obstacle; Pen.\ = maximum penetration depth; Clear.\ = mean minimum
clearance (negative = inside an obstacle on average); Time = mean planning time per
pair. Both planners connect the same measurement poses, so the wrench data and
hence the density estimate (Table~\ref{tab:main}) are unchanged; what differs is
whether the path is executable without contact. Joint-limit violations are zero for both planners by construction
(Sec.~\ref{sec:posepath}). The 3-link condition numbers are structural, not
numerical: one configuration observes total mass and first moment (rank $4$)
against five unknowns, the rank-four ceiling of Prop.~\ref{prop:rank} appearing
directly in the conditioning; the two three-unknown objects are full rank.}}
\end{table}



How many gravity directions per round: two orthogonal directions reach the
rank ceiling, and the orthonormal triad is the knee of diminishing
returns---$12.88$ nat for one direction, $25.41$ for two, $26.35$ for the
triad, only $27.73$ for six---while one direction is anisotropic
($2.72\times$): a gravity-aligned joint axis hides distal mass.
% [번역] 라운드당 중력 방향 개수: 직교 2방향이면 rank 천장에 닿고 직교정규
%   삼면체가 수확체감의 꺾이는 지점이다 — 1방향 12.88 nat, 2방향 25.41,
%   삼면체 26.35, 6방향이라야 27.73 — 단일 방향은 비등방적이다(배향 따라
%   2.72배): 중력과 나란한 관절 축이 그 아래 질량을 숨긴다.

The triad's free $SO(3)$ orientation (Cor.~\ref{cor:invariance}) is spent on
the joints: rotating it to minimize peak torque on those joints preserves
the information \emph{exactly}---$26.3535$ nat before and
after, a $0.0\%$ loss---while cutting peak joint torque from $1.031$ to
$0.620\,\mathrm{N\cdot m}$, a $39.9\%$ reduction: the quantitative basis for
claiming measurement does not endanger the object's joints.
% [번역] 삼면체의 공짜 SO(3) 방향(따름정리 1)은 관절 보호에 쓴다. 그 관절의
%   최대 토크를 최소화하도록 돌리면 정보는 '정확히' 보존되고 — 전후 26.3535
%   nat, 손실 0.0 % — 최대 관절 토크는 1.031 → 0.620 N·m, 39.9 % 감소다.
%   측정이 물체의 관절을 상하게 하지 않는다는 주장의 정량적 근거다.

We grid each joint range at five values, plan every pose pair (triad cycles,
home transits), and score straight-line interpolation against RRT-Connect
(Table~\ref{tab:path}, Fig.~\ref{fig:path}); RRT-Connect never collides, at
negligible path-length and planning-time cost.
% [번역] 관절당 다섯 값 격자의 모든 자세쌍(삼면체 순환, home transit)을
%   계획해 직선 보간 대 RRT-Connect 를 채점한다(표 IV, 그림 6). RRT-Connect
%   는 무충돌 — 경로·계획시간 비용은 무시할 만하다.

The worst \emph{table} penetration is
$33.1\,\mathrm{mm}$ (3-link tip link), while the larger $93.8$ and
$83.5\,\mathrm{mm}$ figures are the object penetrating the robot's own wrist
and forearm. Joint-limit violations are zero for \emph{both} planners, necessarily (the
limits form a convex box): interpolation fails by collision, never by limit
violation. Nor always---the 3-link transits are straight-line
collision-free---so feasibility is re-checked per configuration.
% [번역] 최악 '테이블' 관통은 33.1 mm(3링크 끝 링크), 더 큰 93.8 / 83.5 mm
%   는 물체가 로봇 손목·전완에 파고드는 관통이다. 관절한계 위반은 '두'
%   계획기 모두 0 이며 필연이다(한계는 볼록 상자): 직선 보간의 실패 방식은
%   충돌이지 한계 위반이 아니고, 항상 실패하지도 않는다 — 3링크 transit 은
%   직선도 무충돌이라 실현가능성은 형상마다 재검사한다.



\subsection{Sim-Ready Asset Validation}
\label{sec:asset}

\begin{figure*}[!t]
\centering
\includegraphics[width=0.74\textwidth]{fig_asset.png}
\caption{Asset dynamic fidelity (\emph{simulation-only surrogate}): tip
displacement in three scenarios against the ground-truth-density reference
plant. \rev{[Sim-vs-real version awaits the hardware sessions.]}}
% [번역] asset 동적 충실도(시뮬 전용 대리 검증): 세 시나리오의 끝단 변위를
%   참값 밀도 기준 plant 와 비교. sim-vs-real 판은 하드웨어 세션 대기.
%   ("밀도만 다름"과 2링크 release 둔감성은 IV-E 산문에.)
% [메모] Fig. 7 — 시뮬 대리 검증 그림(figures/v2/fig_asset.png).  구 Fig. 8
%   (실물 비교 자리표시자)은 이 캡션의 rev 로 흡수, fig:assetreal 라벨 보존.
\label{fig:asset}
\label{fig:assetreal}
\end{figure*}

% Table 5 — sim-ready asset 검증.  출처: figures/v2/asset_validation.json (E5).
%   ★ 시뮬-시뮬 대리 검증: reference 는 GT-밀도 plant 이지 현실이 아니다.
%   캡션·헤더에 명시.  8-페이지 압축: 9행 -> push 3행 (Fig. 7 시나리오와 일치),
%   release·drop 순서는 본문 한 문장, 원행은 CUTS.md.  2link release 구조적
%   둔감 서술은 본문 산문에 잔존.  개입/시간 블록: 가이드라인 요구.
%   실물 타이밍은 \Pending.
\begin{table}[!htb]
\caption{Sim-ready asset validation, a \emph{simulation-only surrogate}:
median worst-part trajectory RMSE [mm] against the ground-truth-density
reference plant, \textsc{push} scenario of Fig.~\ref{fig:asset}, with the
user intervention each asset requires.}
% [번역] sim-ready asset 검증(시뮬 전용 대리 검증): 참값 밀도 기준 plant 대비
%   최악 part 궤적 RMSE [mm] 중앙값, 그림 7의 push 시나리오, 아래 블록은
%   asset 별 사용자 개입. (공유 상수·초기조건 5개·release/drop 요약은 각주와
%   IV-E 산문으로 이동.)
\label{tab:asset}
\centering
\scriptsize
\setlength{\tabcolsep}{3pt}
\begin{tabular}{lrrrr}
\toprule
& \multicolumn{4}{c}{\shortstack{Worst-part RMSE [mm], \textsc{push}\\vs.\ GT-density plant (simulation)}} \\
\cmidrule(lr){2-5}
Object & Ours & Uniform & SiPhy & PUGS \\
\midrule
2-link     & \textbf{26.97} & 79.94  & 338.68 & 342.36 \\
3-link     & \textbf{43.97} & 216.13 & 491.38 & 582.44 \\
Stand lamp & \textbf{25.66} & 53.16  & 325.22 & 363.97 \\
\midrule
\multicolumn{4}{l}{\emph{User intervention}} & Time \\
\multicolumn{1}{l}{\quad Ours}
 & \multicolumn{3}{l}{as RORA $+$ scale $+$ joint adj.\ (1--2$\times$)} & \Pending \\
\multicolumn{1}{l}{\quad Uniform}
 & \multicolumn{3}{l}{as RORA $+$ scale; no interaction} & \Pending \\
\multicolumn{1}{l}{\quad SiPhy}
 & \multicolumn{3}{l}{single RGB image; no interaction} & \Pending \\
\multicolumn{1}{l}{\quad PUGS}
 & \multicolumn{3}{l}{multi-view scan; no interaction} & \Pending \\
\bottomrule
\end{tabular}
\vspace{2pt}
\parbox{\linewidth}{\footnotesize \textit{The reference is a plant carrying the
ground-truth densities, \emph{not} the real object; contact stiffness,
dissipation, friction, and joint damping are shared constants across all
variants, so only the per-part densities differ. Medians are over five
pre-registered initial conditions; release and drop are summarized in the text
(Sec.~\ref{sec:asset}). \rev{End-to-end wall-clock time requires the hardware
runs and is pending.} Vision-baseline densities carry the
relative errors of Table~\ref{tab:main} with an assumed positive sign; PhysX-Omni
is excluded because it is N/A on both custom objects.}}
\end{table}



This validation is a \emph{simulation-only surrogate}: the reference is a
Drake plant carrying the ground-truth densities, not the real object, and
nothing here is a real-world claim. \rev{The sim-to-real comparison
(Fig.~\ref{fig:assetreal}) requires the hardware sessions and is pending.}
Contact, friction, and damping parameters are shared constants, so the
comparison isolates exactly what we identify: the per-part densities.
% [번역] 이 검증은 '시뮬레이션 전용 대리 검증' 이다. 기준은 참값 밀도를
%   실은 Drake plant 이지 실제 물체가 아니고, 여기의 어떤 것도 실세계 주장이
%   아니다. sim-to-real 비교(그림 7 내 rev)는 하드웨어 세션 대기로 pending.
%   접촉·마찰·댐핑 계수는 공유 상수 — 우리 식별 대상이 아니다 — 라 비교는
%   부위별 밀도만 격리한다.

Each variant runs release, drop, and push; ours ranks first in all $27$
scenario cells. Table~\ref{tab:asset} reports the push rows matching
Fig.~\ref{fig:asset}; release and drop preserve the same ordering at larger
margins---orders of magnitude, versus $2.1$--$14.2\times$ on push, where
contact amplifies every variant's error, ours included: a ranking, not
absolute accuracy.
% [번역] 각 asset 변형을 release·drop·push 로 구동해 최악 part 궤적 RMSE 로
%   판정하며, 27개 시나리오 칸 전부에서 우리가 1위다. 표 V 는 그림 7과 같은
%   push 행만 싣고, release·drop 은 같은 순서를 더 큰 격차로 보존한다 — 수
%   자릿수 대, 접촉이 모든 변형의 오차를 증폭하는 push 의 2.1~14.2배: push 는
%   절대 정확도가 아니라 순위다.

The 2-link release must not be over-read: it is structurally blind to
density---with the root link welded, the single pendulum's motion is
invariant to scaling the moving link's density, so its sub-millimetre RMSE
reflects that invariance, not skill---and release trajectories
settle onto joint limits, under-stating differences.
% [번역] 2링크 release 는 과대 해석 금지. 구조적으로 밀도에 눈멀어 있다 —
%   뿌리 링크를 용접하면 단진자라 움직이는 링크의 밀도 배율에 불변, sub-mm
%   수치는 실력이 아니라 그 불변성이다 — release 는 관절 한계에 내려앉아
%   차이를 과소평가한다. (vision 변형의 양의 부호 가정은 표 V 각주.)

Intervention is part of the asset's specification (Table~\ref{tab:asset}):
the addition over the inherited RORA~\cite{lee2026rora} stage is the joint
adjustment, once per round---one to two rounds here \rev{(operator and
wall-clock times await the hardware sessions)}.
% [번역] 개입은 asset 사양의 일부다(표 V). 상속되는 RORA 단계 위에 추가되는
%   것은 joint 조정뿐, 라운드당 한 번 — 여기서는 한두 라운드다(작업자·전체
%   소요 시간은 하드웨어 세션 대기).


